% GENERATED FILE. Edit canonical lessons, then run npm run generate:books.
% canonical-source-hash: fnv1a64:0ff40eb2dcbb54e4
% canonical-lessons: BN-C330-utso, BN-C330-srishti-kora, BN-C330-probhab, BN-C330-protikriya, BN-C330-prerona

\chapter{Where It Came From}
\label{ch:bn-where-it-came-from}
\hlchaptermodality{330}

\begin{chapteropening}
\textbf{By the end of this chapter:} \emph{I can say where something started and what it led to.}

Complete the last lesson of chapter 330: I can say where something started and what it led to.
\end{chapteropening}

\section[\texorpdfstring{\bn{উৎস} (utsô)}{উৎস (utsô)}]{\bn{উৎস} (utsô) — a source; an origin}

\label{lesson:BN-C330-utso}



\begin{quote}
Before the new one: say the Bengali for a female singer, then the Bengali for a female dancer.
\end{quote}

\subsection*{You'll want to know: \bn{উৎস}}
\textbf{\bn{উৎস}} — \emph{utsô} — \textquotedblleft{}a source; an origin\textquotedblright{}.

\begin{quote}
\bn{এই} \bn{খবরের} \bn{উৎস} \bn{কী}?
\end{quote}

\begin{quote}
\emph{ei khôbôrer utsô kī?}
\end{quote}

\begin{quote}
Where did this news come from?
\end{quote}

\begin{cousinweb}
\textbf{\bn{উৎস}} (\emph{utsô}) is Sanskrit \emph{utsa}, a spring where water rises (Monier-Williams). Bengali uses it for the place anything first flows from, news or money or a river; that wider use is a later extension of the picture.
\end{cousinweb}

\begin{grammarlens}[title={What you've built}]
One more word for giving reasons and opinions.
\end{grammarlens}

\subsection*{Your turn}
\begin{itemize}
\raggedright
  \item \emph{Say it:} \emph{utsô}
  \item \emph{Say it:} \emph{utsô}, once more
  \item \emph{Say it:} \emph{ei khôbôrer utsô kī?} — Where did this news come from?
  \item \emph{Recall:} say the Bengali for a female singer, then the Bengali for a female dancer, then say \emph{utsô} again
\end{itemize}

\begin{tcolorbox}[breakable,colback=teal!4,colframe=teal!35!black,title={Before you move on}]
What does \bn{উৎস} mean? (\textquotedblleft{}A source; an origin\textquotedblright{}.) Say it once more.
\end{tcolorbox}
\section[\texorpdfstring{\bn{সৃষ্টি} \bn{করা} (srishṭi kôrā)}{সৃষ্টি করা (srishṭi kôrā)}]{\bn{সৃষ্টি} \bn{করা} (srishṭi kôrā) — to create; to give rise to}

\label{lesson:BN-C330-srishti-kora}



\begin{quote}
Before the new one: say the Bengali for a female dancer, then the Bengali for a source.
\end{quote}

\subsection*{You'll want to know: \bn{সৃষ্টি} \bn{করা}}
\textbf{\bn{সৃষ্টি} \bn{করা}} — \emph{srishṭi kôrā} — \textquotedblleft{}to create; to give rise to\textquotedblright{}.

\begin{quote}
\bn{নতুন} \bn{নিয়মটা} \bn{অনেক} \bn{ঝামেলা} \bn{সৃষ্টি} \bn{করল।}
\end{quote}

\begin{quote}
\emph{nôtun niyômṭā ônek jhāmelā srishṭi kôrlô.}
\end{quote}

\begin{quote}
The new rule caused a lot of trouble.
\end{quote}

Past (he/she …): \bn{তিনি} \bn{একটা} \bn{নতুন} \bn{দল} \bn{সৃষ্টি} \bn{করলেন}, \emph{tini ekṭā nôtun dôl srishṭi kôrlen}.

\begin{cousinweb}
\textbf{\bn{সৃষ্টি} \bn{করা}} (\emph{srishṭi kôrā}) is Sanskrit \emph{sṛṣṭi}, creation or bringing forth (Monier-Williams), with \bn{করা} 'to do'. It can make something good (a new team) or something unwelcome (trouble).
\end{cousinweb}

\begin{grammarlens}[title={What you've built}]
One more word for giving reasons and opinions.
\end{grammarlens}

\subsection*{Your turn}
\begin{itemize}
\raggedright
  \item \emph{Say it:} \emph{srishṭi kôrā}
  \item \emph{Say it:} \emph{srishṭi kôrā}, once more
  \item \emph{Say it:} \emph{nôtun niyômṭā ônek jhāmelā srishṭi kôrlô.} — The new rule caused a lot of trouble.
  \item \emph{Recall:} say the Bengali for a female dancer, then the Bengali for a source, then say \emph{srishṭi kôrā} again
\end{itemize}

\begin{tcolorbox}[breakable,colback=teal!4,colframe=teal!35!black,title={Before you move on}]
What does \bn{সৃষ্টি} \bn{করা} mean? (\textquotedblleft{}To create; to give rise to\textquotedblright{}.) Say it once more.
\end{tcolorbox}
\section[\texorpdfstring{\bn{প্রভাব} (prôbhāb)}{প্রভাব (prôbhāb)}]{\bn{প্রভাব} (prôbhāb) — an influence; an effect}

\label{lesson:BN-C330-probhab}



\begin{quote}
Before the new one: say the Bengali for a source, then the Bengali for to create.
\end{quote}

\subsection*{You'll want to know: \bn{প্রভাব}}
\textbf{\bn{প্রভাব}} — \emph{prôbhāb} — \textquotedblleft{}an influence; an effect\textquotedblright{}.

\begin{quote}
\bn{আমার} \bn{ওপর} \bn{এই} \bn{বইটার} \bn{অনেক} \bn{প্রভাব} \bn{আছে।}
\end{quote}

\begin{quote}
\emph{āmār opôr ei bôiṭār ônek prôbhāb āchhe.}
\end{quote}

\begin{quote}
This book has had a big influence on me.
\end{quote}

\begin{cousinweb}
\textbf{\bn{প্রভাব}} (\emph{prôbhāb}) is Sanskrit \emph{prabhāva}, might or power (Monier-Williams). The Bengali sense, the pull one thing has on another, is a softening of that power; the influence falls on something, so it goes with -\bn{এর} \bn{ওপর}.
\end{cousinweb}

\begin{grammarlens}[title={What you've built}]
One more word for giving reasons and opinions.
\end{grammarlens}

\subsection*{Your turn}
\begin{itemize}
\raggedright
  \item \emph{Say it:} \emph{prôbhāb}
  \item \emph{Say it:} \emph{prôbhāb}, once more
  \item \emph{Say it:} \emph{āmār opôr ei bôiṭār ônek prôbhāb āchhe.} — This book has had a big influence on me.
  \item \emph{Recall:} say the Bengali for a source, then the Bengali for to create, then say \emph{prôbhāb} again
\end{itemize}

\begin{tcolorbox}[breakable,colback=teal!4,colframe=teal!35!black,title={Before you move on}]
What does \bn{প্রভাব} mean? (\textquotedblleft{}An influence; an effect\textquotedblright{}.) Say it once more.
\end{tcolorbox}
\section[\texorpdfstring{\bn{প্রতিক্রিয়া} (prôtikriyā)}{প্রতিক্রিয়া (prôtikriyā)}]{\bn{প্রতিক্রিয়া} (prôtikriyā) — a reaction; a response}

\label{lesson:BN-C330-protikriya}



\begin{quote}
Before the new one: say the Bengali for to create, then the Bengali for an influence.
\end{quote}

\subsection*{You'll want to know: \bn{প্রতিক্রিয়া}}
\textbf{\bn{প্রতিক্রিয়া}} — \emph{prôtikriyā} — \textquotedblleft{}a reaction; a response\textquotedblright{}.

\begin{quote}
\bn{খবরটা} \bn{শুনে} \bn{তার} \bn{প্রতিক্রিয়া} \bn{কী} \bn{ছিল}?
\end{quote}

\begin{quote}
\emph{khôbôrṭā shune tār prôtikriyā kī chhilô?}
\end{quote}

\begin{quote}
How did he react when he heard the news?
\end{quote}

\begin{cousinweb}
\textbf{\bn{প্রতিক্রিয়া}} (\emph{prôtikriyā}) is Sanskrit \emph{pratikriyā}, from \emph{prati} 'against, back' and \emph{kriyā} 'an action': an action given back (Monier-Williams gives requital). The everyday sense 'a reaction' is a modern use of the same picture.
\end{cousinweb}

\begin{grammarlens}[title={What you've built}]
One more word for giving reasons and opinions.
\end{grammarlens}

\subsection*{Your turn}
\begin{itemize}
\raggedright
  \item \emph{Say it:} \emph{prôtikriyā}
  \item \emph{Say it:} \emph{prôtikriyā}, once more
  \item \emph{Say it:} \emph{khôbôrṭā shune tār prôtikriyā kī chhilô?} — How did he react when he heard the news?
  \item \emph{Recall:} say the Bengali for to create, then the Bengali for an influence, then say \emph{prôtikriyā} again
\end{itemize}

\begin{tcolorbox}[breakable,colback=teal!4,colframe=teal!35!black,title={Before you move on}]
What does \bn{প্রতিক্রিয়া} mean? (\textquotedblleft{}A reaction; a response\textquotedblright{}.) Say it once more.
\end{tcolorbox}
\section[\texorpdfstring{\bn{প্রেরণা} (preronā)}{প্রেরণা (preronā)}]{\bn{প্রেরণা} (preronā) — an inspiration; a motivation}

\label{lesson:BN-C330-prerona}



\begin{quote}
Before the new one: say the Bengali for an influence, then the Bengali for a reaction.
\end{quote}

\subsection*{You'll want to know: \bn{প্রেরণা}}
\textbf{\bn{প্রেরণা}} — \emph{preronā} — \textquotedblleft{}an inspiration; a motivation\textquotedblright{}.

\begin{quote}
\bn{আমার} \bn{বাবাই} \bn{আমার} \bn{প্রেরণা।}
\end{quote}

\begin{quote}
\emph{āmār bābāi āmār preronā.}
\end{quote}

\begin{quote}
My father is my inspiration.
\end{quote}

\begin{cousinweb}
\textbf{\bn{প্রেরণা}} (\emph{preronā}) is Sanskrit \emph{preraṇa}, setting in motion or urging on (Monier-Williams), in the feminine form. In Bengali it is the person or idea that pushes you forward.
\end{cousinweb}

\begin{grammarlens}[title={What you've built}]
One more word for giving reasons and opinions.
\end{grammarlens}

\subsection*{Your turn}
\begin{itemize}
\raggedright
  \item \emph{Say it:} \emph{preronā}
  \item \emph{Say it:} \emph{preronā}, once more
  \item \emph{Say it:} \emph{āmār bābāi āmār preronā.} — My father is my inspiration.
  \item \emph{Recall:} say the Bengali for an influence, then the Bengali for a reaction, then say \emph{preronā} again
\end{itemize}

\begin{tcolorbox}[breakable,colback=teal!4,colframe=teal!35!black,title={Before you move on}]
What does \bn{প্রেরণা} mean? (\textquotedblleft{}An inspiration; a motivation\textquotedblright{}.) Say it once more.
\end{tcolorbox}
